%% kapa: exam=1 round=29 year=2018 period=2 form=A subject=appraisal_law subject_ko=감정평가관계법규 seq=39 number=39
%% source: Q-Net 감정평가사 기출문제 (cst00309) artlSeq=5202808 file=extracted/1008/5202808/2018년 제29회 감정평가사 1차시험 기출문제/2018년 제29회 감정평가사 1차시험 기출문제/2교시 A형.pdf
%% conversion: kordoc 4.24.51 -> tools/gen1.py (auto); review=proofread-sonnet; flags=none
\begin{question}{39}
도시 및 주거환경정비법령상 비용부담 등에 관한 설명으로 옳지 않은 것은?
\begin{choices}
  \item 정비사업비는「도시 및 주거환경정비법」또는 다른 법령에 특별한 규정이 있는 경우를 제외하고는 사업시행자가 부담한다.
  \item 사업시행자는 토지등소유자로부터 정비사업비용과 정비사업의 시행과정에서 발생한 수입의 차액을 부과금으로 부과ㆍ징수할 수 있다.
  \item 시장ㆍ군수가 아닌 사업시행자는 부과금 또는 연체료를 체납하는 자가 있는 때에는 시장ㆍ군수에게 그 부과ㆍ징수를 위탁할 수 있다.
  \item 국가는 시장ㆍ군수가 아닌 사업시행자가 시행하는 정비사업에 소요되는 비용의 일부에 대해 융자를 알선할 수 없다.
  \item 정비구역안의 국ㆍ공유재산은 정비사업외의 목적으로 매각하거나 양도할 수 없다.
\end{choices}
\end{question}
\answer{4}
